% Theory section (draft) for LPA mechanism
% Will be inserted into adibench_paper.tex as Section "Why does LPA help? An empirical theory"
\section{Why does LPA help? An empirical theory}
\label{sec:theory}

The LPA prototype aggregation (Eq.~(3)) replaces the simple
mean of support features with a weighted mixture that includes
related dialects. We argued informally that this provides a
useful inductive bias. Here we make the argument quantitative.

\paragraph{When does linguistic similarity correlate with feature geometry?}
The mechanism assumes that \emph{linguistically related dialects
have similar AraBERT features}. If this is true, then borrowing
prototypes from related classes effectively borrows more training
signal than noise. Figure~\ref{fig:corr} shows the relationship
between feature-space distances (Euclidean distance between
class-mean AraBERT features) and our dialectological distance
matrix on NADI 18 (18 classes, 200 tweets per class).

\begin{figure}[t]
\centering
\includegraphics[width=0.95\linewidth]{figs/fig5_distance_correlation.pdf}
\caption{\textbf{Distance correlation.} (a) Pairwise feature
distances between class centroids in AraBERT space.
(b) Pairwise linguistic distances derived from
\citet{Versteegh2006,Habash2010}. Spearman $\rho{=}0.435$,
$p < 10^{-15}$: the two matrices are significantly correlated,
meaning \emph{linguistically related dialects \emph{are}
similar in the encoder's feature space}.}
\label{fig:corr}
\end{figure}

\paragraph{Implication.}
The high correlation ($\rho{=}0.435$) confirms that LPA's
similarity-weighted prototype aggregation is doing what it is
designed to do: \emph{leaning on related classes' prototypes
amounts to a sensible regularisation}. If the correlation were
zero (i.e.\ linguistic distance unrelated to feature geometry),
LPA would simply average prototypes of unrelated classes and
add noise; in that case LPA would degrade to vanilla ProtoNet.
Our empirical correlation is strong enough to justify LPA's
\emph{a priori} design and explains why LPA helps.

\paragraph{When does LPA help most?}
LPA is most useful when (i) the encoder's feature geometry
\emph{is} aligned with dialectology, and (ii) the support set
per class is small (few-shot regime), so the prototype estimate
itself is noisy. Empirically (Table~\ref{tab:main}), LPA's
largest gain over ProtoNet comes on amgadhasan 1-shot
($+0.093$, Table~\ref{tab:main}, $+9.3$ percentage points):
the amgadhasan corpus is the smallest (147K tweets) and the
1-shot setting gives the noisiest prototype estimates, which
is exactly when a strong inductive bias pays off.